\section{Privacy and Security Analysis}

\noindent\textbf{Setup.}
$D=(x_1,\ldots,x_n)\in[0,1]^n$,
semi-honest static corruption,
at most 1 party among $\{P_1,P_2,P_3\}$.

\subsection{Differential Privacy Guarantee}

\begin{theorem}[$\varepsilon$-DP of ANCRE]
The mechanism $\mathcal{M}(D) =
\lfloor\hat{\mu}(D)\cdot S\rceil + \eta$,
$\eta\sim\operatorname{DLap}(\Delta_1^S/\varepsilon)$,
$\varepsilon=0.5$, $\delta=0$,
satisfies $\varepsilon$-differential privacy
under substitute adjacency.
\end{theorem}

\begin{proof}[Proof sketch]
TMoM sensitivity:
$\Delta_1 \leq 2/n_{\mathrm{eff}}$.
Standard discrete Laplace mechanism
analysis~\cite{DworkRoth2014} gives
$\varepsilon$-DP. By post-processing
invariance, MPC output preserves DP.
\end{proof}

\subsection{MPC Security}

Under H6-MPC (semi-honest, static,
$\leq 1$ corrupted party):

\begin{itemize}
\item \textbf{$P_1$ corrupted:}
IND-CPA encryption prevents signal recovery.
\item \textbf{$P_2$ corrupted:}
Enclave zeroization — only $\tilde{\mu}$ exits.
\item \textbf{$P_3$ corrupted:}
One Shamir share insufficient for reconstruction.
\end{itemize}

A single corrupted party achieves at most
Denial of Service, not privacy violation.

\subsection{Threat Model Limitations}

We explicitly acknowledge the following
limitations of the current model:

\begin{itemize}
\item \textbf{Semi-honest only:}
Full malicious security (UC framework)
is left for future work.
\item \textbf{Static corruption:}
Adaptive adversaries are not modeled.
\item \textbf{Composition:}
Privacy budget under continuous streaming
requires RDP analysis (future work).
\item \textbf{Secure erasure:}
We assume application-level destruction;
hardware-backed erasure is not formalized.
\end{itemize}
